% =====================================================================
%  四、符号说明
% =====================================================================
\section{符号说明}
本文使用的主要符号见表~\ref{tab:symbol}，未列出的符号在首次出现处说明。

\begin{longtable}{lll}
    \caption{主要符号说明}\label{tab:symbol}\\
        \toprule
        符号 & 含义 & 单位 \\
        \midrule\endfirsthead
        \multicolumn{3}{c}{表~\thetable（续）：主要符号说明}\\
        \toprule 符号 & 含义 & 单位 \\
        \midrule\endhead
        $S_i=(x_i,y_i),\ \theta_i$ & 第 $i$ 个检测点的坐标与示向度 & m, $^\circ$ \\
        $T=(x_T,y_T),\ T_c$ & 单源及频道 $c$ 干扰源的真实位置 & m \\
        $\Delta(\alpha,\beta)$ & 两方位角之差，归一到 $[-180^\circ,180^\circ)$ & $^\circ$ \\
        $\varepsilon,\ \eta$ & 示向度误差界及含舍入余量的角楔半宽 & $^\circ$ \\
        $\mathcal{R},\ V$   & 定位区域及其顶点集 & 无 \\
        $D$                 & 定位区域直径 & m \\
        $M$                 & 直径圆圆心（直径中点） & m \\
        $A,B$               & 直径端点 & m \\
        $S_1,\ S_2$   & 局部测点：$S_1=(0,0),\ S_2=(a,b)$ & m \\
        $\mathcal A_2$      & 第二个检测点的候选区域 & 无 \\
        $x_{max}$           & 干扰源最大可被有效接收距离（$1000\sim1500$ m） & m \\
        $E[D](a,b)$       & 定位区域直径的期望（第二个检测点的价值函数） & m \\
        $\Omega$            & 干扰源所在的靶区 & 无 \\
        $\mathcal C,\ c$             & 频道集合及频道索引 & 无 \\
        $\rho_c$            & 频道 $c$ 对应干扰源的接收半径 & m \\
        $P_c,\ \widetilde P_c$ & 频道 $c$ 的相容定位区域及其保守外包 & 无 \\
        $\hat x_c,\ r_c$     & 源位置估计中心及外包多边形的包围半径 & m \\
        $\mathcal C_{\rm unheard}$ & 尚未发现干扰源的频道集合 & 无 \\
        $h,\ \delta_h$       & 认证方格边长及半对角线 $h/\sqrt2$ & m \\
        $\gamma,\ d(\gamma)$ & 定向源的发射朝向角及单位方向向量 & rad, 无 \\
        $I_s(q),\ Q_c^-$     & 格中心 $q$ 的可排除朝向弧及频道负观测站集 & 无 \\
        $N_{\rm clear,ok}$   & 成功清除次数（亦为成功清除源数） & 无 \\
        $L$                 & 机器人完成任务的总行程 & m \\
        $T_{\rm exit},\ T_{\rm last}$ & 算法退出时刻与末次成功清除时刻 & s \\
        \bottomrule
\end{longtable}
